% Section "imitation" (integration of 2026-10-14, research/paper-review/imitation-edit/log-integrate.md). Compact form of the
% draft in research/tracks/imitation/notes-final.md §2.5, cut as in its §5.3; corrections C1-C6 of
% research/tracks/imitation/referee-final.md §6 applied (C1-C5 here and in app-imit.tex, C6 in discussion.tex); literature
% sentences of research/paper-review/imitation-edit/literature.md §3.2 (I1 in the first paragraph; I2 in the last paragraph,
% condensed because the related-work paragraph carries the full attributions). Placement: after sec:limits, before
% sec:computable. Sources: research/tracks/imitation/robust/notes-final.md (Lemma 2.1, Thm 2.2, Props 3.1, 3.3, Setting 4.1,
% Lemmas 4.4-4.6, Thm 4.7(U),(vii), Cor 4.9, Props 4.10(a2),(c),(d), 4.11, Prop 5.2) with referee-A.md (A1, R1),
% referee-B.md (B1-B3) and imitation/referee-final.md (items 1-5, all verified); imitation/notes-final.md §2.2 (B_beta
% conditioned on (h_0, mu) only). Moved to app-imit.tex: rem:imit:private (app:imit:private), the numbers after the theorem
% (app:imit:main), prop:imit:limits(c) and the cell count (now prop:imit:cheap, app:imit:cheap), and the in-context learner
% (Bayes mixture over persons; robust Prop 5.2) to app:imit:private. q:imit:open is dropped; its
% questions are the open-problems item "Imitation" in discussion.tex. Local notation: Sigma (coin set; B is the truncation
% level of sec:teacher), rho (coin law; not a structure function), beta (coupling), W (cells), bar mu, P^*.
% Fix round of 2026-10-14 (research/paper-review/imitation-edit/log-fix.md): imitation-edit review-math M1, N2 (worlds renamed), N3; review-front M1 (paragraphs 1-2), m2, m3.
\section{Imitation: a sampler whose coins the environment drives}\label{sec:imit}

Imitation motivates both reasons: a learner trained on a person's answers, as in behavioural cloning \citep{bain1999framework,dehaan2019causal}, should return a copy of the person. In our terms the person is a sampler, a model with fresh coins at each query, and a \emph{copy} is a program whose kernel equals the person's on the support of the inputs. If the person's coins are private, this is the first reason: no predictor gains over the person more than its length plus a Ville deficiency, however it models the environment (\cref{rem:imit:private}), and every fixed $\lambda>1$ eventually returns a shortest copy, not a fill-in (\cref{thm:recovery:main}). Systems that infer from the data which person they face are excluded by the choice of the private-randomness class, not by the data (\cref{app:imit:private}).

Here the person's coins depend on the environment, so hypotheses that model the environment can predict the person's answers \emph{better} than the person does. The data follow the kernel $P^*$ of \cref{def:imit:setting}, so a fixed $\lambda$, or any $\lambda_n=o(n)$, eventually returns a shortest program for $P^*$ (\cref{thm:recovery:main}(a)), not a copy when $P^*$ differs from the person's kernel; in the feature case such a program encodes the coupling. The coins are reused at every input of a cell, like an error rule, and the guarantee below needs the schedule of the second reason, $\lambda/n$ above the reuse factor $2^{-m}$ (\cref{sec:teacher}). When the coins are an incompressible function of a variable, the cell, with many more values than the person has bits, we show that for $\lambda/n$ in a window every near-optimal private-randomness hypothesis, whatever it computes, is close to the person in behaviour (\cref{thm:imit:main}), and that no procedure reliably copies each of several persons who share one data law (\cref{prop:imit:limits}). We do not show that a copy itself is selected or that the selected program computes as the person does, and a coupling with a short description defeats every $\lambda/n$ (\cref{prop:imit:cheap}); the guarantee is for the population objective, from samples only asymptotically. Proofs are in \cref{app:imit}.

\begin{definition}[Environment-driven coins]\src{imitation/robust notes-final Setting 4.1 (B$_\beta$ conditioned on $(h_0,\mu)$ only, imitation/notes-final \S2.2; imitation/referee-final item 3); referee-B B1}\label{def:imit:setting}
Let $Y$ and the coin set $\Sigma$ be finite, $m\ge1$, and $W:=\{0,1\}^m$, the \emph{cells}.
(a) A \emph{channel} $\mu=(\mu(\cdot\mid w))_{w\in W}$ of finitely supported laws on inputs with rational masses gives the input law $\bar\mu:=2^{-m}\sum_w\mu(\cdot\mid w)$, with support $\supp\mu$; in the \emph{feature case} every input in $\supp\mu(\cdot\mid w)$ has $w$ as its first $m$ bits.
(b) The \emph{person} has a coin law $\rho$ on $\Sigma$, every $\rho(b)2^m$ a positive integer, and a response function $h:\supp\mu\times\Sigma\to Y$, injective in $b$; its kernel is $\nu_H(y\mid x):=\rho(b)$ if $y=h(x,b)$, and $0$ otherwise, and its \emph{model} is $h_0:=(h,\rho)$. A \emph{copy} is a program of $\Hcl{\times}$ whose kernel equals $\nu_H$ on $\supp\mu$; $\mathcal E$ is the set of copies, $\ell_0:=\min_{\mathcal E}\ell$, and $\pi_H$ is a shortest copy.
(c) A \emph{coupling} $\beta:W\to\Sigma$ of exact type $\rho$, $|\beta^{-1}(b)|=\rho(b)2^m$, drives the coins: at each query, independently, $w_i$ is uniform on $W$, $x_i\sim\mu(\cdot\mid w_i)$ and $y_i:=h(x_i,\beta(w_i))$. The learner sees $(x_i,y_i)$, i.i.d.\ from $\bar\mu\otimes P^*$ with $P^*(y\mid x):=\sum_w\mu(x\mid w)\ind[y=h(x,\beta(w))]/\sum_w\mu(x\mid w)$.
(d) $R(\pi):=\Ex[-\log\nu_\pi(y_i\mid x_i)]$, $J_c:=c\,\ell+R$ and $\Phi(c):=\inf_{\Hcl{\times}}J_c$; every copy has $R=H(\rho)$. $L_\beta:=\log\big(2^m!/\prod_b(\rho(b)2^m)!\big)$, $B_\beta:=K(\beta\mid h_0,\mu)$, $\delta_\beta:=L_\beta-B_\beta$, and $I_\beta(\pi):=B_\beta-K(\beta\mid h_0,\pi,\mu)$, the information of $\pi$ about the coupling. For a probability $P$ and a sub-probability $Q$, $D_{1/2}(P\Vert Q):=-2\log(1-\Hel(P,Q))$.
\end{definition}

The coin used at a query is a fixed function of the cell, which the learner sees in the input (the feature case) or only through the channel (a latent shared cause). The data identify $P^*$ (\cref{thm:recovery:main}); when the input carries information about the coin $\beta(w)$, that is, when $P^*\ne\nu_H$ on a set of positive $\bar\mu$-mass, hypotheses that model $\beta$ predict better than the person, by up to $H(\rho)$ bits per sample in the feature case, where $P^*$ is deterministic. The teacher of \cref{def:teacher:instance} is the feature case with $h(x,b)=f_a(x)\oplus b$, $\rho=\Bern(p)$ and $\beta=\ind_S$: systematic errors are environment-driven coins. As for error rules, $L_\beta\ge2^mH(\rho)-|\Sigma|\log(2^m+1)$ \citep{cover2006elements}, and $\Prob[\delta_\beta>j]<2^{-j}$ for a coupling drawn uniformly from its type class, independently of the person's model and the channel.

\begin{theorem}[Robust near-imitation]\status{proved}\src{imitation/robust notes-final Thm 4.7(U),(vii), Lemmas 4.4--4.6; referee-A A1; imitation/notes-final \S2.2 (re-derived); imitation/referee-final items 1--3 (verified; code q1--q4); checks f7, s1(B),(C)}\label{thm:imit:main}
In the setting of \cref{def:imit:setting} let $c=a2^{-m}$ with $a>1$, $\epsilon\ge0$, and $G:=\ell_0+\delta_\beta+|\Sigma|\log(2^m+1)+\kappa$. Every $\pi\in\Hcl{\times}$ with $J_c(\pi)\le\Phi(c)+\epsilon$, whatever its internal structure, satisfies:
\begin{enumerate}[label=(\roman*),itemsep=0pt,topsep=2pt]
\item $\ell(\pi)\le\ell_0+(G+2^m\epsilon)/(a-1)$;
\item $-(aG+2^m\epsilon)/((a-1)2^m)\le R(\pi)-H(\rho)\le c\,\ell_0+\epsilon$: it predicts neither much better nor much worse than the person;
\item $I_\beta(\pi)\le\ell_0+(G+2^m\epsilon)/(a-1)+\kappa$;
\item with $E:=2\log\big(2^{m-1}(2H(\rho)+c\,\ell_0+\epsilon)+1\big)\le2m+2\log(2H(\rho)+c\,\ell_0+\epsilon+1)$,
\[
\Ex_{x\sim\bar\mu}D_{1/2}\big(\nu_H(\cdot\mid x)\,\Vert\,\nu_\pi(\cdot\mid x)\big)\ \le\ c\,\ell_0+\frac{2a}{a-1}\cdot\frac{G+E+\kappa}{2^m}+\frac{a+1}{a-1}\,\epsilon,
\]
so its Hellinger gap from the person, $\Ex_{x\sim\bar\mu}\Hel(\nu_H(\cdot\mid x),\nu_\pi(\cdot\mid x))$, is at most $\frac{\ln2}2$ times the right side.
\end{enumerate}
Moreover $J_c(\pi_H)\le\Phi(c)+G/2^m+c\,\ell_0$. In the feature case, if $|\Sigma|\ge2$ and the values of $\rho$ are distinct, a program of length at most $\ell_0+\lceil L_\beta\rceil+a_\rho+\kappa$, $a_\rho:=2\log(m+1)+2\log(|\Sigma|+1)+|\Sigma|(m+1)$, has kernel $P^*$ and risk $0$, and so beats every copy for $c<H(\rho)/(\lceil L_\beta\rceil+a_\rho+\kappa)$, a threshold at least $2^{-m}(1+(a_\rho+1+\kappa)/(2^mH(\rho)))^{-1}$.
\end{theorem}

\emph{Proof idea.} A program is a code for the coupling, $K(\beta\mid h_0,\pi,\mu)\le2^mR(\pi)+\kappa$ (\cref{lem:imit:coding}), which gives the whole class a right line of slope $2^{-m}$ at $(\ell_0,H(\rho))$; \cref{lem:kink:main} gives (i), (ii) and the bound on $\Phi$. For (iv), the divergence is at most that of per-cell geometric means of the predictions (\cref{lem:imit:closeness}), which define a semimeasure on couplings; normalised, it would give $\beta$ a code shorter than the coupling's incompressibility allows, unless the divergence is small, the program predicts worse than the person, or it knows more about $\beta$.

The best predictor has $\Ex_xD_{1/2}(\nu_H\Vert P^*)=H(\rho)$ in the feature case. So when the coupling is incompressible ($\delta_\beta\ll2^mH(\rho)$) and the person much simpler than it ($\ell_0+|\Sigma|m\ll2^mH(\rho)$), for every $c$ above the reuse factor $2^{-m}$, with $a$ bounded away from $1$, and well below $H(\rho)/\ell_0$, every near-optimal hypothesis is close to the person in behaviour, on the scale $H(\rho)$, and holds little of the coupling, although hypotheses that model the environment predict better. The theorem gives near-copies, not copies: nothing excludes a nearby coin law as the minimiser (cf.\ \cref{rem:teacher:selected}); and the hypotheses that predict better gain only by storing an incompressible table. By \cref{prop:recovery:beyond}(c) with the sampler $P^*$, almost surely every $cn$-MAP satisfies (i)--(iv) with $\epsilon=0$ for all large $n$. What happens outside this range of $c$, and numerical illustrations, are in \cref{app:imit:main}.

\begin{proposition}[What no procedure can do]\status{proved; computed}\src{imitation/robust notes-final Props 4.10(a2), 4.11; referee-A R1; imitation/referee-final item 4 (verified; code q5); checks f8, s1(D)}\label{prop:imit:limits}
\begin{enumerate}[label=(\alph*),itemsep=1pt,topsep=2pt]
\item \emph{One data law, at most one copy.} If worlds $\mathcal W_1,\dots,\mathcal W_N$ as in \cref{def:imit:setting} have one law of the data pairs and persons whose kernels differ pairwise on sets of positive $\bar\mu$-mass, then every procedure $\Psi$ (a measurable map from samples and independent randomness to programs) has $\sum_i\Prob_{\mathcal W_i}[\Psi(h_n)\in\mathcal E(\nu_i)]\le1$ at every $n$. In the feature case with $|\Sigma|\ge2$, one such world has the person $P^*$, who responds to the cell.
\item \emph{A person entangled with its coupling.} Let $X=W$, $\mu(\cdot\mid w)=\delta_w$, $Y=\Sigma=\{0,1\}$, $\rho=\Bern(s_0/2^m)$ with $1\le s_0<2^{m-1}$, and let the person $H_f$ answer $f(w)\oplus b$ for a table $f:W\to\{0,1\}$. Draw $f$ uniformly and $\beta$ uniformly from its type class, independently. Then $\Prob[\delta_\beta>j]<2^{-j}$, yet every procedure copies the person with probability at most $2^{-L_\beta}$, at every $n$.
\end{enumerate}
\end{proposition}

\emph{Proof idea.} (a) The output has one law in all the worlds, and the sets of copies are disjoint. (b) The data depend on the world only through $f\oplus\ind_\beta$, which $2^{L_\beta}$ equally likely worlds with pairwise different persons share.

So imitation data identify behaviour, not mechanism: persons with the same response function and coin law share a data law, however they compute them, and the selected program, a shortest one for its own kernel, may compute it unlike the person and extrapolates off $\supp\bar\mu$ as it does. Nor do the data identify the split of $P^*$ into a person and an environment, as a policy does not determine its split into a planner and a reward \citep{armstrong2018occam}: every procedure adopts a convention, the $c$-procedure by rate (\cref{prop:limits:separable}). An incompressible coupling is not enough: in (b) the person is as complex as its coupling and responds to the same cells, so a hypothesis such as $\ell_0\ll2^m$, or the left condition that list decoding gives in \cref{thm:teacher:main}, is needed. A coupling with a short description, such as a pseudo-random generator with a short seed, defeats every $c$ (\cref{prop:imit:cheap}); if it is the person's own deterministic state, the program that models it is arguably the more faithful copy, and the target $\nu_H$ is itself a modelling choice. Data from several environments or from interventions identify more in inverse reinforcement learning and behavioural cloning \citep{cao2021identifiability,dehaan2019causal}; whether such data identify which structure belongs to the person here is open (\cref{sec:discussion:open}).
